\documentclass[10pt,a4paper]{amsart}
\usepackage[margin=19mm]{geometry}
\usepackage{amsmath,amssymb,mathtools}
\usepackage[scr=rsfs,cal=euler]{mathalfa}
\usepackage{xfrac}
\usepackage[unicode,hidelinks,bookmarks=false]{hyperref}
\newcommand{\ie}{i.e.\ }
\newcommand{\cf}{cf.\ }
\newcommand{\resp}{respectively\ }
\newcommand{\Subsection}{\subsection}
\providecommand{\nobreakdash}{\nobreak\hbox{-}}
\setlength{\emergencystretch}{2em}
\multlinegap0pt
\usepackage{amssymb}
\usepackage[shortlabels]{enumitem}
\newtheorem{proposition}{Proposition}
\title{On Diophantine pairs and triples of triangular numbers}
\author{Marija Bliznac Trebje\v{s}anin}
\date{Source: arXiv:2603.29565v2 (CC BY 4.0)}
\begin{document}
\maketitle

\noindent\emph{Source and licence.} On Diophantine pairs and triples of triangular numbers. Version arXiv:2603.29565v2 (CC BY 4.0), distributed by arXiv under the Creative Commons Attribution 4.0 International licence, http://creativecommons.org/licenses/by/4.0/, as recorded in the arXiv licence metadata for this exact version and shown on the abstract page https://arxiv.org/abs/2603.29565v2. Full licence notice: \texttt{licenses/arxiv-2603.29565-CC-BY-4.0.txt}.\par\smallskip

\noindent\textit{Editorial scope.} Counted unit: the article's proposition prop:nonexist (source0.tex lines 180--183). The article's complete original proof of it is delivered with it (source0.tex lines 185--223, one \texttt{proof} environment of 39 lines). No mathematics is written, repaired or re-derived: the statement and the proof are the article's own lines, in the article's own notation, and no retained line is substituted or re-flowed.  Retained uncounted as the source-local context the proof needs: lines 57--59.  Nothing was omitted from any retained span, so the package carries no omission record.  No retained line contains a citation, so the wrapper carries no bibliography and no bibliography entry.  No retained line contains a figure environment, an image-inclusion command or drawing code, so no image is delivered and the package declares no asset dependency.  Editorial formatting only, in addition: an AMS \texttt{amsart} preamble that restates the article's own declarations and macros used by the retained lines, namely the environments \texttt{proposition} on separate counters, together with the standard packages those lines need.  The retained environments print in this document as Proposition 1 in order of appearance, whereas the article prints the same items with the numbering of its own full document.  The complete native source of the article as distributed in the arXiv e-print for this exact version is bundled unchanged as \texttt{sources/197523/source0.tex}.
\begin{equation}\label{pellova_x}
    x^2-n(n+1)y^2=16a-n(n+1),
\end{equation}

\begin{proposition} \label{prop:nonexist}
Let $a \in \{12, -37\}$. For any positive integer $n$, the equation \eqref{pellova_x}
does not have any integer solutions $(x, y)$ such that $y$ is odd. Consequently, there exists no pair of triangular numbers $\{T_n, T_m\}$ with the property $D(12)$ or $D(-37)$.
\end{proposition}

\begin{proof}
Assume, for the sake of contradiction, that there exists a positive integer $n$ such that equation \eqref{pellova_x} has an integer solution $(x, y)$ with an odd integer $y \ge 1$. 
By setting $y = 2m+1$ for some positive integer $m$, equation \eqref{pellova_x} can be rewritten as:
$$ x^2 = 16a + n(n+1)((2m+1)^2-1). $$
Notice that $n(n+1) = \frac{N_0^2-1}{4}$. Hence, it is useful to introduce the notation $N_0 = 2n+1$, $M_0 = 2m+1$ and $c = 2|x|$, multiply the equation by $4$ and obtain the following symmetric formulation:
\begin{equation} \label{eq:sym_combined_mixed}
c^2 = (N_0^2-1)(M_0^2-1) + 64a.
\end{equation}
By symmetry, we may assume without loss of generality that $m \geq n$, which implies $M_0 \geq N_0$. In what follows, we restrict our analysis of equation \eqref{eq:sym_combined_mixed} to this case.

If $M_0 = 1$ (or $N_0 = 1$), then  $c^2 = 64a$. 
 For $a=12$, this yields $c^2 = 768$, which is not a perfect square.
    For $a=-37$, this yields $c^2 = -2368$, which is trivially impossible.

Hence, any valid solution to \eqref{eq:sym_combined_mixed} requires $N_0 \ge 3$ and $M_0 \ge 3$. 

Equation \eqref{eq:sym_combined_mixed} can be rewritten as a generalized Pellian equation in variables $X=c$ and $Y=M_0$:
$$ X^2 - (N_0^2-1)Y^2 = 64a - (N_0^2-1). $$
If it admits any solutions with $M_0 \ge N_0 \ge 3$, then by the well-ordering principle, there must exist a solution for a fixed $N_0$ such that the positive integer $M_0$ is minimal. 
Let $(c, M_0)$ be this minimal integer solution. The fundamental solution $(N_0, 1)$ to the associated Pell equation $X^2 - (N_0^2-1)Y^2 = 1$ allows us to define a new solution $(c', M_0')$:
$$ (c', M_0') := \bigl(N_0 c - (N_0^2-1)M_0,\;\; N_0 M_0 - c\bigr). $$
Observe that the parity of the new elements $(c', M_0')$ is preserved, and that $(c', |M_0'|)$ is also a solution to equation \eqref{pellova_x}. 
Furthermore, $|M_0'| \neq 1$; otherwise, we would have either $16a=(m-n)^2$ or $16a=(m+n+1)^2$, which occurs only when $a$ is a perfect square. Since $M_0'$ is odd, we conclude that $|M_0'| \ge 3$.

For the new solution $(c', |M_0'|)$ to satisfy the  inequality $|M_0'| < M_0$, the value of $c$ must fall strictly between $(N_0-1)M_0$ and $(N_0+1)M_0$. Squaring this condition yields two inequalities that must hold simultaneously:
\begin{align}
    c^2 &< (N_0+1)^2 M_0^2 \implies 2M_0^2(N_0+1) > 64a + 1 - N_0^2, \label{upper_bound} \\
    c^2 &> (N_0-1)^2 M_0^2 \implies 2M_0^2(N_0-1) > -64a - 1 + N_0^2. \label{lower_bound}
\end{align}
If both bounds hold, we obtain a new integer solution with $3 \le |M_0'| < M_0$, which contradicts our assumption that $M_0$ is minimal. Therefore,  at least one of these inequalities must fail. This restricts the minimal solution to a finite set of exceptional pairs $(N_0, M_0)$.


\textbf{Case $a=12$:} For a positive parameter, \eqref{lower_bound} holds trivially. Inequality \eqref{upper_bound} fails when $2M_0^2(N_0+1) \le 769 - N_0^2$. This yields exactly six pairs with $M_0 \ge N_0 \ge 3$: $(3,3)$, $(3,5)$, $(3,7)$, $(3,9)$, $(5,5)$, and $(5,7)$. Substituting these into \eqref{eq:sym_combined_mixed} generates values for $c^2$ which are not perfect squares.


 \textbf{Case $a=-37$:} For a negative parameter, inequality \eqref{upper_bound} holds trivially. The lower bound does not hold when $2M_0^2(N_0-1) \le 2367 + N_0^2$. From equation \eqref{eq:sym_combined_mixed}, the condition $c^2 \ge 0$ imply that $(N_0^2 - 1) (M_0^2 - 1) \geq 2368$. Combining these two inequalities yields 13 pairs: $(3,19)$, $(3,21)$, $(3,23)$, $(5,11)$, $(5,13)$, $(5,15)$, $(5,17)$, $(7,9)$, $(7,11)$, $(7,13)$, $(9,9)$, $(9,11)$, and $(11,11)$. However, none of these pairs yields a perfect square for $c^2$.

Since the supposed minimal solution $(c, M_0)$ must belong to this set of pairs, and none of these pairs yield a perfect square for $c^2$, such a minimal solution cannot exist. Consequently, the equation has no solutions for $a=12$ or $a=-37$.
\end{proof}

\end{document}
